\gddchapter{High-Level Overview}


\section{Game Concept}


It's year 2030, post-apocalipse. Your brother never returned from a scavenger mission in a nearby shopping mall.
You decide to rescue him yourself no matter the price.

\section{Genre and Platform}
Text based rougelike with a core mission.

\section{Platform}
The game is intended to be run on PC via a web browser.


\section{Audience}

The audience is mostly people who enjoy imaginary experiences such as ttrpg games
and other games that employ deep immersive states.

\section{2 variations of the game}

The game constists of two version necessary for the research for the thesis that's linked to the game.

\begin{enumerate}
    \item The keyboard version - In this version the player chooses their actions from a predetermined
    list of possibilities, this way they can navigate with certainty 
    \item The voice version - This version of the game employs voice commands and the player doesn't see 
    all of the possibilities in front of them. Instead of that they rely on the description of their current 
    environment and come up with the actions themselves
\end{enumerate}

The rationale behind having that voice interation is that the player may potentially deepen his immersion
due to him having an illusion of greater control of the game while also formulating the actions with his voice
might also lead to increased immersion compared to using manual button input. The validity of these claims
will be tested by the acoompaning research.